\chapter{Divisibilité et nombres premiers}
\label{corriges:03}
\section*{Exercice 3.1 -- Décomposer et simplifier}
$420=2^2\times3\times5\times7$ et $630=2\times3^2\times5\times7$. On simplifie les facteurs communs $2\times3\times5\times7=210$ :
\[\frac{420}{630}=\frac23.\]
Les entiers $2$ et $3$ n'ont aucun diviseur commun supérieur à $1$.

\section*{Exercice 3.2 -- Lots identiques}
$96=2^5\times3$ et $144=2^4\times3^2$. Le plus grand diviseur commun est $2^4\times3=48$. On obtient donc $48$ lots de $96/48=2$ biscuits et $144/48=3$ chocolats. Un nombre de lots plus grand ne pourrait pas diviser les deux quantités.

\section*{Exercice 3.3 -- Deux signaux périodiques}
$45=3^2\times5$, $60=2^2\times3\times5$. Le plus petit multiple commun positif vaut $2^2\times3^2\times5=180$ secondes, soit $3$ minutes. La prochaine coïncidence est à 14 h 03.

\section*{Exercice 3.4 -- Premier ou composé ?}
$221=13\times17$ est composé. Pour $223$, $\sqrt{223}<15$, donc les seuls diviseurs premiers à tester sont $2,3,5,7,11,13$. Le nombre est impair, sa somme de chiffres vaut $7$, il ne finit ni par $0$ ni par $5$, et
$223=7\times31+6=11\times20+3=13\times17+2$.
Aucun test n'aboutit : $223$ est premier.
